\documentclass[aspectratio=169,11pt]{beamer}
% For Chinese slides: \documentclass[aspectratio=169,11pt]{ctexbeamer}, compiled with XeLaTeX.

\usepackage{amsmath,amssymb}
% Add \usepackage{graphicx} or pgfplots only when the deck includes figures; call figures by relative path, e.g. \includegraphics[width=\linewidth]{figures/fig1.png} (no \graphicspath).

\setbeamertemplate{navigation symbols}{}
\setbeamertemplate{headline}{}
\setbeamertemplate{footline}{%
  \begin{beamercolorbox}[wd=\paperwidth,ht=2.25ex,dp=1ex,center]{page number in head/foot}%
    \usebeamerfont{page number in head/foot}\insertframenumber
  \end{beamercolorbox}%
}

\title{Paper Title}
\subtitle{Author (Year), Journal}
\author{Presenter}
\date{}

\begin{document}

\begin{frame}
  \titlepage
\end{frame}

\begin{frame}{Introduction}
  \begin{itemize}
    \item The question the paper asks
    \item Why it matters
    \item The answer, previewed
  \end{itemize}
\end{frame}

% Titles are "Topic: point"; the topic is the same on every slide of a section.
\begin{frame}[label=m-core]{The Model: The Main Result}
  \textbf{Proposition 1.} The paper's statement of the result, in its own words:
  \[
    \alert{Y} = F(X_1, X_2)
  \]
  \begin{itemize}
    \item Define each symbol where it first appears
    \item Why the result holds
  \end{itemize}
  \vfill\hfill\hyperlink{a-derivation}{\beamergotobutton{Appendix: derivation}}
\end{frame}

\begin{frame}{The Model: Why It Holds}
  \begin{itemize}
    \item The next step of the same part, under the same topic
    \item One sentence of intuition
  \end{itemize}
\end{frame}

\begin{frame}{Conclusion}
  \begin{itemize}
    \item One line per takeaway
  \end{itemize}
\end{frame}

\appendix

\begin{frame}[label=a-derivation]{Appendix}
  \framesubtitle{Derivation of Proposition 1}
  Secondary derivations and long algebra go here; steps of the core argument stay in the main slides.

  \vfill\hfill\hyperlink{m-core}{\beamerreturnbutton{Back}}
\end{frame}

\end{document}
